\chapter{Materiales e Instrumental}
\begin{itemize}
    \item \textbf{Multímetro Digital (Patrón Experimento 1):} Pro'sKit MT-1210 (S/N: 15224613). Exactitud: $\pm (1 \percent \text{ lect.} + 2 \text{ dig.})$
        (hasta $200 \V$). Ver su hoja de datos en el \textbf{Anexo \ref{anexo:hoja_datos}}.
    \item \textbf{Multímetro Analógico (A contrastar):} Electrónico UNIVO. Exactitud declarada: $\pm 2,5 \percent$ del
        Fondo de Escala.
    \item \textbf{Multímetros Digitales (Experimento 2):} UNI-T Modelo UT33A+ (2 unidades, empleados simultáneamente como voltímetro de precisión en rango $2\V$ y miliamperímetro en rango $200\mA$).
    \item \textbf{Fuente de alimentación:} Regulable de $0 \V$ a $30 \V$.
    \item \textbf{Circuito generador de corriente constante:} Módulo basado en la fuente SK150 (en reemplazo del
        regulador LM317).
    \item \textbf{Probeta a ensayar:} Tramo de cable arrollado de forma no inductiva (plegado por la mitad y luego
        arrollado, longitud $12\text{ m} \pm 12\text{ cm}$).
    \item \textbf{Telurímetro:} UNI-T Modelo UT522.
\end{itemize}